\section*{Task 4: Fish Farming}

\begin{enumerate}[label=\alph*)]
    \item\label{item:task4a} The population of fish $P(t)$ in a pond can be described by the following ODE
    \begin{equation}\label{eq:fish}
    \frac{d}{dt}P(t) = rP(t)\left(1 - \frac{P(t)}{C}\right)
    \end{equation}
    The constants $r$ and $C$ specify the growth rate and maximum capacity of the pond. \\\\
    \textbf{Subtask:} \\
    Use IFDIFF to solve the ODE on the time interval $T = [0, 5]$ with initial conditions $P(0) = 0.1$. \\
    Set the constants as $r = 2$ and $C = 1$ directly in the RHS (do not pass them as parameters). \\
    Plot the solution. What is the maximum size of the population? When does the population grow fastest? \\ \\
    \textbf{Steps:}
    \begin{enumerate}[label=\arabic*.]
        \item Implement the RHS in the file \texttt{rhsTask4.m} as described in \cref{eq:fish}. Remember to set the constants $r$ and $C$ directly in the RHS.
        \item Set the integration time span and initial conditions in the file \texttt{runTask4.m}. Ignore the parameters for now.
        \item Use IFDIFF to prepare the RHS for integration. Don't forget to include the integrator and its options.
        \item Use IFDIFF to solve the initial value problem. You can then plot your solution using the provided \lstinline!plotPopulation(solution)! function.
    \end{enumerate}
    \textbf{Hints:}
    \begin{itemize}
        \item Check the appendix for a reminder on how to formulate and solve an ODE using IFDIFF.
    \end{itemize}

    \item The population of fish will now be harvested automatically according to the following rule:\\
    Whenever the population $P(t)$ would exceed a threshold $T \in [0, C]$, a proportion of the current population will be extracted from the pond and added to the harvest $H(t)$ according to the harvest ratio $\alpha \in [0, 1]$.

    This new system can be described by the following ODE
    \begin{equation}\label{eq:fish_harvest}
        \frac{d}{dt}\begin{bmatrix}P(t) \\ H(t)\end{bmatrix} = \begin{bmatrix}rP(t)\left(1 - \frac{P(t)}{C}\right) \\ 0\end{bmatrix}
    \end{equation}
    with state jumps
    \begin{equation}\label{eq:fish_jump}
        \begin{bmatrix}P(t^+) \\ H(t^+)\end{bmatrix} = \begin{bmatrix}P(t^-) \\ H(t^-)\end{bmatrix} +  \begin{bmatrix}-\alpha P(t^-) \\ \phantom{-}\alpha P(t^-)\end{bmatrix}, \quad P(t^-) = T
    \end{equation}

    \textbf{Subtask:} \\
    Use IFDIFF to solve the new system for the same constants and initial conditions as in subtask \ref{item:task4a} and $H(0) = 0$. \\
    Make sure to pass $p = (T, \alpha)$ as parameters to your RHS (do not set them directly in the RHS). \\
    Plot the solution. What is the maximum size of the population now? Play around with different parameters. \\ \\
    \textbf{Steps:}
    \begin{enumerate}[label=\arabic*.]
        \item Extend your RHS from subtask \ref{item:task4a} in the file \texttt{rhsTask4.m}. Split the state vector into two components and compute the derivatives as described in \cref{eq:fish_harvest}. \\
        Use the functions \lstinline|ifdiff_jumpif(condition, direction)| and \lstinline|ifdiff_update(jump)| to implement the state jump as described in \cref{eq:fish_jump}. \\
        Make sure that the components of your state and parameter vectors are defined in the same way as in the exercise sheet, i.e.~$y = (P, H)^T$ and $p = (T, \alpha)$.
        \item Use your code from subtask \ref{item:task4a} to solve the new system and plot the solution. Try out different values for $T$ and $\alpha$.
    \end{enumerate}
    \textbf{Hints:}
    \begin{itemize}
        \item Check the appendix for a reminder on how to formulate state jumps in a RHS using IFDIFF.
        \item Verify that your states and parameters are defined correctly: $y = (P, H)^T$ and $p = (T, \alpha)$.
    \end{itemize}

     \item We now want to examine how the solution behaves depending on the parameters $T$ and $\alpha$. \\ \\
     \textbf{Subtask:} \\
     Discretize the parameter space $p = (T, \alpha) \in [0, C] \times [0, 1]$ as a 2D-grid. \\
     Use IFDIFF to compute a solution for each parameter in the grid and plot the results. \\
     Examine the behavior at the grid boundaries, i.e.~for $T = 0,\: T = C$ and $\alpha = 0, \:\alpha = 1$.
\end{enumerate}
